\setcounter{section}{0}
\begin{center}
	\zihao{3}
	\textbf{河北工程大学2022$\sim*$2023学年第二学期期末考试(A)卷}
\end{center}

\section{单项选择题}

\begin{enumerate}
	\item 若级数 $\sum_{n=1}^{\infty}u_n$ 收敛于 $S$，则 $\sum_{n=1}^{\infty}u_{n+1}$ 必\underline{\hspace{2cm}}。
	\begin{choices}
		\item 收敛于 $S$
		\item 收敛于 $S+u_1$
		\item 收敛于 $S-u_1$
		\item 发散
	\end{choices}
	
	\item 三重积分
	$
	\iiint_{x^2+y^2+z^2\le 1}\frac{x\ln(x^2+y^2+z^2+1)}{\sqrt{x^2+y^2+z^2+1}}\,dx\,dy\,dz=\underline{\hspace{2cm}}.
	$
	\begin{choices}
		\item $0$
		\item $\pi$
		\item $2\pi$
		\item $4\pi$
	\end{choices}
	
	\item 设 $L$ 为单位圆周 $x^2+y^2=1$ 的右半部分，则
	$
	\int_L e^{x^2+y^2}\,ds=\underline{\hspace{2cm}}.
	$
	\begin{choices}
		\item $\pi$
		\item $2e$
		\item $\pi e$
		\item $2\pi e$
	\end{choices}
	
	\item 函数$z = f(x,y)$在点 $(x,y)$ 可微是 $f(x,y)$ 在该点偏导数存在的\underline{\hspace{2cm}}条件。
	\begin{choices}
		\item 充分必要
		\item 必要
		\item 充分
		\item 既非充分又非必要
	\end{choices}
	
	\item 设 $f(x)$ 是周期为 $2\pi$ 的周期函数，在 $[-\pi,\pi)$ 上的表达式为 $f(x)=x+1$，则 $f(x)$ 的傅里叶级数在点 $x=\pi$ 处收敛于\underline{\hspace{2cm}}。
	\begin{choices}
		\item $-\pi+1$
		\item $\pi+1$
		\item $0$
		\item $1$
	\end{choices}
\end{enumerate}

\section{填空题}

\begin{enumerate}
	\item 交换二次积分
	$
	\int_0^1 dx\int_0^{\sqrt{x}} f(x,y)\,dy
	$
	的次序为\underline{\hspace{2cm}}。
	
	\item 若区域 $D:x^2+y^2\le 1,\ x\ge 0,\ y\ge 0$，则二重积分
	$
	\iint_D \sqrt{1-x^2-y^2}\,dx\,dy=\underline{\hspace{2cm}}.
	$
	
	\item 设 $L$ 为圆心在原点的单位圆（取逆时针方向），则
	$
	\oint_L (x+y)\,dx+(x-y)\,dy=\underline{\hspace{2cm}}.
	$
	
	\item 若 $z=f(x^2+y^2)$，且 $f$ 一阶可导，则
	$
	\frac{\partial z}{\partial x}=\underline{\hspace{2cm}}.
	$
	
	\item 方程
	$
	xy''+2x^2y'^4+x^3y=x^4+1
	$
	是\underline{\hspace{2cm}}阶微分方程。
\end{enumerate}

\section{计算题}

\begin{enumerate}
	\item 求
	$
	z=x^3-y^3+3x^2+3y^2-9x
	$
	的极值。
	
	\item 计算二重积分
	$
	\iint_D xy\,dx\,dy,
	$
	其中 $D$ 为由直线 $y=1,\ x=2$ 及 $y=x$ 所围成的区域。
	
	\item 设 $L$ 是曲线 $y^2=x$ 上从 $A(1,-1)$ 到点 $B(1,1)$ 的有向弧段，计算
	$
	\int_L x\,y\,dx.
	$
	
	\item 求曲线
	$
	x=t^2,\quad y=1-t,\quad z=t^3
	$
	在点 $(1,0,1)$ 处的切线方程和法平面方程。
	
	\item 判定级数
	$
	\sum_{n=1}^{\infty}(-1)^n\frac{n}{3^n}
	$
	的敛散性，如果收敛，是绝对收敛还是条件收敛。
	
	\item 已知平面曲线通过点 $(0,1)$，且该曲线上任一点 $P(x,y)$ 处切线的斜率为 $4x^3y$，求该曲线的方程。
\end{enumerate}

\section{解答题}

\begin{enumerate}
\item 求微分方程
$
y''-5y'+6y=e^{5x}
$
的通解。


\item 计算曲面积分
$
\iint_{\Sigma}(z^2+x)\,dy\,dz-zx\,dx\,dy,
$
其中 $\Sigma$ 为旋转抛物面
$
z=\frac{x^2+y^2}{2}
$
介于平面 $z=0$ 及 $z=2$ 之间的部分取下侧。



\item 求级数
$
\sum_{n=1}^{\infty}\frac{x^n}{n}
$
的收敛域及和函数，并求数项级数
$
\sum_{n=1}^{\infty}\frac{1}{n\cdot 2^n}
$
的和。

\item 设
$
f(x,y)=
\begin{cases}
	(x^2+y^2)\sin\frac{1}{x^2+y^2}, & x^2+y^2\ne 0,\\
	0, & x^2+y^2=0.
\end{cases}
$
试判断 $f(x,y)$ 在 $(0,0)$ 点处是否连续，偏导数是否存在，是否可微。

\end{enumerate}
\newpage
